% Table 3 -- operational definitions of the evaluation modes and the two contrast tests.
% These resolve the boundary cases a benchmark survey must adjudicate (mirrors P1 Table 3).
\begin{table*}[t]
\centering
\footnotesize
\setlength{\tabcolsep}{5pt}
\renewcommand{\arraystretch}{1.3}
\resizebox{\textwidth}{!}{%
\begin{tabular}{@{}p{0.19\textwidth} p{0.40\textwidth} p{0.30\textwidth}@{}}
\toprule
\textbf{Construct} & \textbf{Counts (a benchmark has it if and only if\ldots)} & \textbf{Does not count} \\
\midrule
\textbf{Open-loop eval (\dopen)} &
it scores prediction or generation \emph{quality} on fixed inputs without executing the
prediction in a feedback loop: video fidelity (FVD), physics or rollout consistency, or
frame-level accuracy &
any protocol where an action is executed and the environment responds; offline scoring of a
policy's success \\
\textbf{Closed-loop eval (\dclosed)} &
success is measured by running a policy \emph{in} the environment with state feedback and
scoring task outcome: success rate, sub-goal completion, or reward under interaction &
scoring a generated video or a static prediction; a metric computed without stepping the
environment \\
\textbf{Bridge (\dbridge)} &
a prediction (imagined rollout, video, or latent plan) is \emph{converted into executed actions}
and scored on the resulting task success &
a benchmark that stops at prediction quality; one that never turns a prediction into a control
signal \\
\textbf{Native VLA-vs-WM contrast ($\gamma$)} &
the benchmark \emph{itself} runs both a direct Vision-Language-Action policy and a world-model
policy under one protocol and reports the head-to-head &
a model-agnostic suite that merely \emph{hosts} either family; a suite evaluated with only one
family in the original paper \\
\textbf{Advantage-of-prediction ($\delta$)} &
a metric quantifies the closed-loop \emph{gain} of a predictive policy over a matched direct-VLA
baseline on the same tasks (assessed, not adopted) &
absolute success with no matched baseline; open-loop quality scores; a ranking that never
isolates the predictive component \\
\bottomrule
\end{tabular}}
\caption{Operational definitions used to place each benchmark on the evaluation axes
(\S\ref{sec:method}). They resolve the ambiguous boundary cases: hosting a policy is not the same
as \emph{building} a VLA-vs-world-model contrast ($\gamma$), and an absolute success number is not
an advantage-of-prediction measurement ($\delta$) without a matched baseline. \dopen\ open-loop,
\dclosed\ closed-loop, \dbridge\ bridge.}
\label{tab:defs}
\end{table*}
